\documentclass{article}
\usepackage[utf8]{inputenc}
\usepackage[spanish]{babel}
\usepackage{graphicx}

\begin{document}
%Portada
  \begin{figure}[lt]
    \includegraphics[width=0.3\textwidth]{ull.eps}
  \end{figure}
  \title{Integración Del Trapecio} %[En el intervalo [1, 6] de una función]
  \author{Tiffany López Nicholson \\ Miriam Martín Jacinto \\ Sergio Vega García}
  \date{\today}
  \maketitle

  \begin{abstract}
    \begin{center}
       A continuación se presentará como se ha implementado con python un algoritmo capaz de resolver la integral definida $\int_{1}^{6} \frac{1}{1+e^x}$
    \end{center}
  \end{abstract}
  \pagebreak

%%%%%%%%%%%%%%%%%%%%%

  \tableofcontents
  \pagebreak
  
%%%%%%%%%%%%%%%%%%%%%
  
  \section{Motivaciones y objetivos}
    \begin{itemize}
      \item Objetivo principal: Implementar con Python la integración del trapecio para f(x)
      \item Objetivo específico: Realizar un algoritmo capaz de resolver la integración del trapecio para $ f(x) = \frac{1}{1 + e^{x}}$, x E [1, 6],
    \end{itemize}
    \subsection{Seccion uno}
    \subsection{Seccion dos}
  \pagebreak

%%%%%%%%%%%%%%%%%%%%%
  \section{Fundamentos teóricos}
    La regla del trapecio es un método de integración numérica. Esta regla se basa en aproximar el valor de la integral de f(x) por el de la función lineal que pasa a través  de ésta es igual al área del trapecio que se forma bajo la gráfica de la función lineal.
    
    \begin{figure}[h]
      \begin{center}
	\includegraphics[scale=0.5]{areadeltrapecio.eps}
	\caption{Ejemplo}
      \end{center}
    \end{figure}
      
    \subsection{Explicación general}

      Para la utilización del método del trapecio partimos de una función, la cual dividiremos en n trozos iguales. Cuanto mayor sea el número de particiones, mayor precisión tendrá el método.\\
      %Intentar introducir aquí un ejemplo gráfico.

      Se puede apreciar que el área tomada por exceso, es decir, la que supera a la función, o la tomada por defecto, la cual no llega a la función, se reduce según aumenta el número de particiones. \\

      La función general es: 
        
      \begin{center}
   
        $\int_{a}^{b}f(x)dx \approx \frac{h}{2}[f(a)-2f(a+h)+2f(a+2h)+...+f(b)]$

      \end{center}
 
      Donde $h=\frac{b-a}{n}$ y n es el número de divisiones.\\

      La expresión anterior también se puede escribir como:

      \begin{center}

        $\int_{a}^{b}f(x)dx \approx \frac{b-a}{n}(\frac{f(a)+f(b)}{2}+\sum\limits_{k=1}^{n-1}f(a+k\frac{b-a}{n}))$
      \end{center}

    \subsection{Explicación aplicada}
      
   \pagebreak

\end{document}
